\title{The Era of 1-bit LLMs: All Large Language Models are in 1.58 Bits}
Additionally, this innovation facilitates novel computational frameworks and catalyzes the development of specialized hardware tailored for 1-bit LLM workloads. In this paper, we propose a novel 1-bit LLM variant, designated as \textbf{\text{BitNet b1.58}}, where each parameter (or weight) assumes ternary values from \{-1, 0, 1\}. To investigate the token scalability of \text{BitNet b1.58}, we trained it on 2T tokens utilizing the same data pipeline as StableLM-3B~\cite{StableLM-3B-4E1T}, a leading open-source model with 3B parameters. Additionally, we reported full inference energy costs for processing sequences of 512 tokens, observing that as model size increases, the energy efficiency advantage of \text{BitNet b1.58} over the FP16-based \text{LLaMA LLM} becomes more substantial. This modification not only streamlines the implementation and benefits system-level efficiency but also has minimal impact on empirical performance. Some strategies have attempted to address this by increasing SRAM capacity to boost throughput, but this approach is substantially more expensive compared to scaling DRAM. \textbf{Native Support for Long Sequences in LLMs}

As LLM applications increasingly require the processing of extended sequences, managing long contexts is a growing necessity. In comparison to their full-precision counterparts, 1-bit LLMs are markedly more efficient regarding memory usage in both storage capacity and bandwidth requirements. Notably, the 3.9B variant of \text{BitNet b1.58} not only surpasses the speed (by 2.4$\times$) and memory efficiency (by 3.32$\times$) of the \text{LLaMA LLM} 3B model, but also demonstrates superior overall performance. The second advantage, confirmed by our empirical results, is that \text{BitNet b1.58} is able to achieve parity with full-precision (FP16) models in both perplexity and downstream task metrics, provided the models share equivalent settings (such as size and the amount of training tokens) at the 3B parameter scale and beyond. \end{itemize}

\paragraph{Training with 2T Tokens}

The volume of training tokens significantly impacts large language model effectiveness. These systems exhibit state-of-the-art results across diverse natural language processing benchmarks; however, their escalating model sizes have introduced obstacles for widespread deployment, especially considering the environmental and financial burdens imposed by their substantial energy usage. While it remains grounded in the foundational BitNet design, it incorporates several notable updates, detailed as follows. \section{Discussion and Future Work}

\textbf{1-bit Mixture-of-Experts (MoE) LLMs}

Mixture-of-Experts (MoE) architectures have emerged as a resource-efficient solution in the context of LLMs. Zero-shot accuracy outcomes are summarized in Table~\ref{tab:2t}. This improvement stems from the scaling behavior of the \emph{nn.Linear} operation, whose time complexity increases with the number of parameters. Mirroring the perplexity analysis, downstream task outcomes confirm that \text{BitNet b1.58} at 3.9B parameters consistently exceeds \text{LLaMA LLM} 3B, while requiring fewer computational resources. In line with this trend, the design of \text{BitNet b1.58} incorporates components inspired by LLaMA. This model is trained from the beginning, utilizing 1.58-bit weights in conjunction with 8-bit activations. Consequently, floating-point additions and multiplications dominate the overall computational cost. Pipeline parallelism~\cite{gpipe} was implemented to enable the execution of \text{LLaMA LLM} 70B within these hardware constraints. \item BitNet b1.58 at 30B parameters exceeds the efficiency of a 7B FP16 LLM across latency, memory, and energy usage. It maintains comparable energy efficiency and offers considerable improvements in memory requirements, throughput, and inference latency over FP16-based LLMs. The results demonstrate that \text{BitNet b1.58} consistently outperforms across all evaluated tasks, underscoring the robust generalization capacity of 1.58-bit LLMs. The results, as depicted in Figure~\ref{fig:latency-memory-energy}, demonstrate that both latency and memory usage trends improve with increased model scale; notably, the performance benefit (in terms of speed-up) becomes more pronounced for larger models. Measurements were obtained via the FasterTransformer\footnote{\url{https://github.com/NVIDIA/FasterTransformer}} framework, which is specifically optimized for accelerating LLM inference on GPU platforms. We assessed zero-shot capabilities over several language benchmarks: ARC-Easy~\cite{arc}, ARC-Challenge~\cite{arc}, Hellaswag~\cite{hellaswag}, Winogrande~\cite{winoGrande}, PIQA~\cite{piqa}, OpenbookQA~\cite{openbookqa}, and BoolQ~\cite{boolq}. The computational approaches pioneered by BitNet~\cite{bitnet} motivate and justify this new hardware direction. Table~\ref{tab:zero-shot} details zero-shot accuracy for a set of downstream tasks. A principal limitation in this scenario is the significant memory usage stemming from KV cache storage during inference. Additionally, \text{BitNet b1.58} introduces two further benefits. Both models underwent evaluation on a composite benchmark comprising Winogrande~\cite{winoGrande}, PIQA~\cite{piqa}, SciQ~\cite{sciq}, LAMBADA~\cite{lambada}, and ARC-easy~\cite{arc}. If the entirety of a model is able to reside on a single chip, this overhead could be eliminated altogether. \section{Results}

We conducted a comparison between \text{BitNet b1.58} and our reimplemented FP16 version of \text{LLaMA LLM} across multiple model scales. By augmenting the original 1-bit BitNet representation with an explicit 0 value, the new design requires just 1.58 bits per parameter in the binary system. This advantage is attributable to the growing share of computation spent on \emph{nn.Linear} operations, while the contribution of other elements becomes relatively less significant in larger models. Specifically, \text{BitNet b1.58} 70B achieves a 4.1-fold reduction in inference time relative to the \text{LLaMA LLM} baseline. The field has advanced from 16-bit implementations towards reduced precision variants, including 4-bit models~\cite{gptq, awq}. To ensure an equitable evaluation, both models underwent pre-training on the RedPajama dataset~\cite{redpajama}, each processing 100 billion tokens. For benchmarks reporting both accuracy and normalized accuracy, we use their average. Further, we analyzed GPU memory usage and inference latency for both \text{LLaMA LLM} and \text{BitNet b1.58}. \paragraph{Energy}

We also calculated the energy expenditure associated with arithmetic operations for both \text{BitNet b1.58} and \text{LLaMA LLM} models. \textbf{LLMs for Edge and Mobile Applications}

Employing 1.58-bit LLMs holds significant promise for enhancing language model functionality on edge and mobile platforms. On a deeper level, the 1.58-bit LLM offers a fresh scaling law and a training strategy that enables future LLM generations to reach both high performance and heightened cost-effectiveness. Furthermore, their compatibility with CPUs—which are the predominant processors in edge and mobile hardware—means that \text{BitNet b1.58} can be effectively run on such devices, advancing their potential capabilities. \paragraph{Throughput}

We assess the throughput performance of \text{BitNet b1.58} relative to the \text{LLaMA LLM} with 70B parameters, deploying both models across two 80GB A100 GPUs. Evaluations followed the \emph{lm-evaluation-harness}\footnote{\url{https://github.com/EleutherAI/lm-evaluation-harness}} procedure. The first is a heightened representational power—attributed to the explicit inclusion of zero in the parameter space, which allows for direct feature filtering and substantively enhances 1-bit LLM effectiveness. To highlight inference efficiency, we reported the average time required per output token. Additionally, the communication cost associated with transferring activations across devices is drastically lowered. \textbf{\text{BitNet b1.58} introduces a new paradigm in the scaling law for model performance vis-à-vis inference cost}. For \text{BitNet b1.58}, the integration included Ladder's~\cite{ladder} 2-bit kernel. Figure~\ref{fig:energy} breaks down the sources of energy consumption: in \text{BitNet b1.58}, most operations involve INT8 addition, whereas \text{LLaMA LLM} predominantly relies on FP16 addition and multiplication. In contrast, BitNet’s matrix multiplication relies solely on integer additions, dramatically reducing the energy consumption typically required by LLMs. With decreased demands on memory and energy, 1.58-bit LLMs make feasible deployments that were previously impractical, unlocking a variety of novel application scenarios. The batch size was incrementally increased up to the maximum allowable by GPU memory, with the input sequence length fixed at 512. It features RMSNorm~\cite{rmsnorm}, SwiGLU~\cite{swiglu}, and rotary embeddings~\cite{rope}, and omits all bias terms. \paragraph{Memory and Latency}

We expanded our experiments to include models with 7B, 13B, and 70B parameters and analyzed the associated computational costs. By quantizing weights and activations, memory and computation requirements are significantly lowered, enabling LLMs to be deployed more efficiently. This model is able to attain parity with its full-precision counterparts (i.e., FP16 or BF16 Transformer LLMs) of equivalent architecture and token count, achieving similar perplexity and downstream task results. \paragraph{Quantization Function.} To ensure the weights are limited to the values -1, 0, or +1, we employ an \emph{absmean} quantization method. Drawing from the analyses depicted in Figure~\ref{fig:latency-memory-energy} and \ref{fig:energy}, we can establish the following correspondences for comparable efficiency between 1.58-bit and 16-bit model configurations:

\begin{itemize}
    \item BitNet b1.58 at 13B parameters surpasses the 3B FP16 LLM in terms of latency, memory footprint, and energy consumption. These decisions facilitate straightforward adoption of \text{BitNet b1.58} into widely used open-source frameworks such as Huggingface, vLLM~\cite{vllm}, and llama.cpp\footnote{\url{https://github.com/ggerganov/llama.cpp}}. Instead, all activations are rescaled per token to fit the range $[-Q_b, Q_b]$, eliminating zero-point quantization. Significantly, starting from the 3B size, \text{BitNet b1.58} can achieve results comparable to the full precision baseline. \paragraph{LLaMA-alike Components.}

LLaMA~\cite{llama, llama2} has established itself as the reference backbone in the open-source LLM landscape. Leveraging 1.58-bit LLMs offers remedies to these bottlenecks. \textbf{Specialized Hardware for 1-bit LLMs}

Initiatives like Groq\footnote{\url{https://groq.com/}} have illustrated the considerable promise of developing dedicated hardware, such as LPUs, to support LLM operations. Despite the reduction in computational FLOPs, practical adoption of MoE is often constrained by high memory demands and the cost of inter-chip communication. Memory utilization displays a comparable pattern, since the embedding layer—maintained in full precision—accounts for a proportionally smaller memory footprint in larger architectures. The accuracy results demonstrate that, as model size increases, the difference between \text{BitNet b1.58} and \text{LLaMA LLM} performance diminishes. The data for both latency and memory were collected using a 2-bit kernel, indicating that further efficiency gains may be achievable through continued optimization. Beyond raw computation, the act of moving model weights from DRAM into an on-chip accelerator’s local memory (such as SRAM) during inference can incur significant overheads. \text{BitNet b1.58} preserves all the core advantages of its predecessor, including its computational framework that nearly eliminates multiplications during matrix operations and can be extensively optimized. In this paper, we propose an advanced 1-bit LLM variant named \textbf{\text{BitNet b1.58}}, in which each parameter adopts one of three possible values: \{-1, 0, 1\}. Our primary consideration was matrix multiplication, given its dominant role in large language model computational cost. The performance statistics for StableLM 3b with 2T tokens were extracted directly from its official documentation. Since the primary bottleneck for computational speed on many hardware platforms is power consumption, these energy reductions also enable acceleration of computation. This technique first normalizes the weight matrix by dividing by its mean absolute value and subsequently rounds each element to the closest value in the set \{-1, 0, +1\}:

\begin{equation}
    \widetilde{W} = \mathrm{RoundClip}(\frac{W}{\gamma + \epsilon}, -1, 1),
\end{equation}

\begin{equation}
    \mathrm{RoundClip}(x, a, b) = \max(a, \min(b, \mathrm{round}(x))),
\end{equation}

\begin{equation}
    \gamma = \frac{1}{nm}\sum_{ij} |W_{ij}|. \section{BitNet b1.58}

\text{BitNet b1.58} builds upon the BitNet framework, a variant of Transformer networks that substitutes \emph{nn.Linear} layers with \emph{BitLinear} layers. \end{abstract}

\section{The Era of 1-bit LLMs}

Artificial Intelligence research has recently witnessed an extraordinary expansion in both the scale and proficiency of Large Language Models (LLMs). In Table~\ref{tab:ppl}, we present a comparative overview of perplexity and resource consumption for both \text{BitNet b1.58} and \text{LLaMA LLM}. These platforms generally operate under stringent memory and compute constraints, which limit both model scale and achievable performance. \end{equation}

For activations, the quantization follows the standard approach in BitNet, with one key distinction: we omit scaling activations to the interval $[0, Q_b]$ before non-linear transformations. A promising avenue to mitigate these challenges is post-training quantization, wherein models are transformed to operate at lower numerical precision after training~\cite{smoothquant, gptq, quip, quip_sharp}. This improvement leads to reduced latency and lower costs when transferring model weights from DRAM, thereby facilitating quicker and more effective inference. Additionally, validation perplexity was measured on the WikiText2~\cite{wikitext2} and C4~\cite{c4} corpora. Primarily, the minimized memory footprint enables MoE models to require fewer hardware units for deployment. Future work could focus on further compressing these activations, potentially to 4 bits or even lower for 1.58-bit LLMs, while preserving information without loss. Nevertheless, post-training quantization, though prevalent in production-level LLMs, is fundamentally limited in optimality. \text{BitNet b1.58} marks an important advancement toward innate support for long sequences, primarily by lowering activation precision from 16 bits to 8 bits, thereby effectively allowing the context window to expand twofold with equivalent computational resources. These results position \text{BitNet b1.58} as a Pareto improvement in the context of state-of-the-art LLMs. Recent advancements in 1-bit model architectures, notably BitNet~\cite{bitnet}, highlight a compelling approach for lowering the operational expenses of LLMs without compromising their accuracy. The findings indicate that \text{BitNet b1.58} achieves parity with full precision \text{LLaMA LLM} in terms of perplexity from the 3B parameter regime onwards, while delivering a 2.71-fold increase in speed and reducing GPU memory utilization by a factor of 3.55. Traditional LLMs employ 16-bit floating-point representations (such as FP16 or BF16), with matrix multiplications constituting the core computational workload. Importantly, this 1.58-bit configuration yields dramatic reductions in latency, memory footprint, computational throughput, and energy demands, thereby delivering substantial improvements in efficiency. \item BitNet b1.58 at 70B parameters is more resource-efficient than a 13B FP16 LLM, measured by latency, memory, and energy metrics. 
\begin{abstract}
Recent advances, exemplified by BitNet~\cite{bitnet}, are ushering in a new generation of 1-bit Large Language Models (LLMs). According to Table~\ref{tab:throughput}, the 70B \text{BitNet b1.58} supports batch sizes up to 11 times larger than those feasible with \text{LLaMA LLM}, translating to an 8.9-fold improvement in throughput. Expanding on these efforts, there is a need to pursue the design and implementation of hardware and systems explicitly tailored for the efficient execution of 1-bit LLMs. Utilizing the energy estimation framework from~\cite{energycost, pokebnn}, we found that \text{BitNet b1.58} reduces the matrix multiplication energy requirement by a factor of 71.4 when run on 7nm semiconductor technology.